\documentclass[a4paper,11pt]{article}
\usepackage[T2A]{fontenc}
\usepackage[utf8]{inputenc}
\usepackage[russian]{babel}
\usepackage{amsmath,amssymb,mathtools}
\renewcommand{\familydefault}{\sfdefault}
\usepackage{icomma}
\usepackage{geometry}
\geometry{left=21mm,right=21mm,top=22mm,bottom=23mm}
\usepackage{microtype}
\usepackage{tikz}
\usetikzlibrary{calc,arrows.meta,positioning,shapes.geometric}
\usepackage{tkz-euclide}
\usepackage{tabularx,booktabs}
\usepackage{enumitem}
\usepackage{setspace}
\usepackage{xcolor}
\usepackage{fancyhdr}
\usepackage{hyperref}
\definecolor{accent}{HTML}{008080}
\definecolor{darkaccent}{HTML}{005959}
\definecolor{textgray}{HTML}{333333}
\color{textgray}
\hypersetup{colorlinks=true,linkcolor=darkaccent,urlcolor=darkaccent}
\onehalfspacing
\setlength{\parindent}{0pt}
\setlength{\parskip}{3pt plus 1pt}
\setlength{\headheight}{23pt}
\setlist[enumerate]{leftmargin=*,itemsep=7pt,topsep=5pt}
\setlist[itemize]{leftmargin=*,itemsep=3pt,topsep=3pt}
\pagestyle{fancy}
\fancyhf{}
\fancyhead[L]{\small\color{textgray}Прямоугольный треугольник и медиана}
\fancyhead[R]{\small\color{textgray}Тимофей Васильченко}
\fancyfoot[C]{\small\thepage}
\renewcommand{\headrulewidth}{0.45pt}
\renewcommand{\headrule}{\hbox to\headwidth{\color{accent}\leaders\hrule height \headrulewidth\hfill}}
\makeatletter
\renewcommand\section{\@startsection{section}{1}{\z@}{2.2ex plus 1ex minus .2ex}{1.1ex plus .2ex}{\normalfont\Large\color{darkaccent}}}
\makeatother
\newcommand{\ex}[1]{\par\medskip\noindent\textcolor{darkaccent}{\textbf{#1}}\par\smallskip}
\newcommand{\key}[1]{\noindent\textcolor{darkaccent}{\textbf{Главная идея.}} #1}
\newcommand{\checkitem}[1]{\item[$\square$] #1}
\begin{document}
\thispagestyle{empty}
\begin{center}
\vspace*{28mm}
{\large\color{darkaccent}МАТЕМАТИКА}\par
\vspace{19mm}
{\LARGE Чек-лист}\par
\vspace{12mm}
{\LARGE\color{accent}Медиана, высота и окружность\\[3pt]
в прямоугольном треугольнике}\par
\vspace{16mm}
{\normalsize Подготовка к ЕГЭ по профильной математике}\par
\vfill
{\large Тимофей Васильченко}\par
\vspace{9mm}
{20 ноября 2026 г.}\par
\vspace{20mm}
{\small Лёвин Артём Александрович}\par
{\small эксклюзивно для Тимофея Васильченко}\par
\vspace{20mm}
\begin{tikzpicture}[x=1cm,y=1cm]
  \draw[accent,line width=0.85pt] (0,0) .. controls (1.4,0.65) and (3.0,-0.35) .. (4.1,0)
                            .. controls (5.6,0.58) and (7.4,-0.42) .. (8.5,0);
\end{tikzpicture}
\end{center}
\clearpage

\section{1. Главная теорема: гипотенуза и диаметр}
\key{В прямоугольном треугольнике центр описанной окружности совпадает с серединой гипотенузы.}

Пусть $\angle ACB=90^\circ$, $M$ --- середина гипотенузы $AB$. Тогда
\[
AB=2R,\qquad MA=MB=MC=R=\frac{AB}{2}.
\]
\begin{center}
\begin{tikzpicture}[scale=.75]
  \tkzDefPoint(0,0){A}
  \tkzDefPoint(6,0){B}
  \tkzDefPoint(3,0){M}
  \tkzDefPoint(1.5,2.598076){C}
  \tkzDrawCircle[color=accent,line width=.8pt](M,A)
  \tkzDrawPolygon[line width=.85pt](A,C,B)
  \tkzDrawSegment[color=darkaccent,dashed](C,M)
  \tkzDrawPoints(A,B,C,M)
  \tkzMarkRightAngle[size=.3](A,C,B)
  \tkzLabelPoint[below left](A){$A$}
  \tkzLabelPoint[below right](B){$B$}
  \tkzLabelPoint[below](M){$M$}
  \tkzLabelPoint[above](C){$C$}
\end{tikzpicture}
\end{center}
\textbf{Почему это верно.} По теореме Фалеса угол, опирающийся на диаметр окружности, прямой. Обратно: описанная окружность прямоугольного треугольника имеет гипотенузу своим диаметром. Поэтому её центр --- середина гипотенузы.

\ex{Пример 1. Радиус описанной окружности}
Если гипотенуза равна $4$, то $R=4/2=\boxed{2}$.

\section{2. Медиана к гипотенузе и обратное свойство}
\begin{itemize}
\item \textbf{Прямое свойство:} медиана, проведённая из вершины прямого угла к гипотенузе, равна её половине: $CM=AB/2$.
\item \textbf{Обратное свойство:} если в треугольнике $ABC$ точка $M$ --- середина $AB$ и $CM=AM=BM$, то $\angle ACB=90^\circ$.
\end{itemize}
\textbf{Обоснование обратного свойства.} Окружность с центром $M$ и радиусом $MA$ проходит через $A$, $B$ и $C$. Отрезок $AB$ является диаметром, значит, угол $ACB$ прямой.

\textbf{Следствие для углов.} Треугольники $AMC$ и $BMC$ равнобедренные:
\[
\angle MAC=\angle ACM,\qquad
\angle MBC=\angle BCM.
\]
Это позволяет находить острые углы исходного треугольника, используя их части при медиане.

\ex{Пример 2. Медиана делит прямой угол в отношении $1:2$}
Пусть $\angle ACB=90^\circ$, $CM=m$, $\angle ACM: \angle MCB=1:2$. Тогда $\angle ACM=30^\circ$, $\angle MCB=60^\circ$. В равнобедренном $AMC$ угол $A=30^\circ$, значит, $\angle B=60^\circ$; гипотенуза $AB=2m$. По свойству треугольника с углами $30^\circ$ и $60^\circ$:
\[
\boxed{BC=m,\qquad AC=m\sqrt3,\qquad AB=2m.}
\]
\clearpage

\section{3. Периметры и теорема Пифагора}
Когда дан периметр, начинайте с формулы $P=a+b+c$ и выразите неизвестные стороны. Если треугольник прямоугольный, обязательно проверьте связь катетов и гипотенузы:
\[
a^2+b^2=c^2.\]
\ex{Пример 3. Периметры треугольников с медианой}
Пусть $\angle ABC=90^\circ$, $M$ --- середина $AC$, $P_{ABM}=8$, $P_{BMC}=9$. Найдём стороны $ABC$.

Обозначим $AC=t>0$, $AB=x>0$, $BC=y>0$. По свойству медианы $AM=BM=CM=t/2$. Поэтому
\[
\begin{cases}
x+t=8,\\ y+t=9,\\ x^2+y^2=t^2.
\end{cases}
\quad\Longrightarrow\quad
\begin{cases}
x=8-t,\\y=9-t,\\(8-t)^2+(9-t)^2=t^2.
\end{cases}
\]
Последнее уравнение приводится к виду $t^2-34t+145=0$, откуда $t=5$ или $t=29$.

\textbf{Отбор корней.} Поскольку $x=8-t>0$, нужно $t<8$; значит, $t=29$ не подходит. При $t=5$ получаем $x=3$, $y=4$. Проверка: $3^2+4^2=5^2$, оба периметра равны заданным.
\[
\boxed{AB=3,\qquad BC=4,\qquad AC=5.}
\]

\section{4. Высота к гипотенузе через две площади}
Пусть $\angle ABC=90^\circ$, $BH\perp AC$. Обозначим $AB=a$, $BC=b$, $AC=c$, $BH=h$. Одну и ту же площадь запишем двумя способами:
\[
S_{ABC}=\frac{ab}{2}=\frac{ch}{2}
\quad\Longrightarrow\quad
\boxed{h=\frac{ab}{c}}.
\]
\begin{center}
\begin{tikzpicture}[scale=1.02]
\tkzDefPoint(0,0){A}
\tkzDefPoint(0,2){B}
\tkzDefPoint(4,2){C}
\tkzDefPoint(0.8,0.4){H}
\tkzDefPoint(2,1){M}
\tkzDrawPolygon[line width=.85pt](A,B,C)
\tkzDrawSegment[color=accent,line width=1.1pt](B,H)
\tkzDrawSegment[color=darkaccent,dashed](B,M)
\tkzMarkRightAngle[size=.23](A,B,C)
\tkzMarkRightAngle[size=.22](B,H,C)
\tkzDrawPoints(A,B,C,H,M)
\tkzLabelPoint[below left](A){$A$}
\tkzLabelPoint[above left](B){$B$}
\tkzLabelPoint[above right](C){$C$}
\tkzLabelPoint[below right](H){$H$}
\tkzLabelPoint[below right](M){$M$}
\end{tikzpicture}
\end{center}
\textbf{Метод:} при известной высоте к одной стороне сравните две формулы площади, выбрав разные основания. Точное значение высоты полезно далее для тригонометрии.

\clearpage
\section{5. Угол между высотой и медианой}
В предыдущем рисунке $M$ --- середина гипотенузы $AC$. Тогда $BM=c/2$, а $\angle BHM=90^\circ$. Для угла $\varphi=\angle HBM$ прямоугольного треугольника $BHM$:
\[
\cos\varphi=\frac{BH}{BM}=\frac{ab/c}{c/2}
=\boxed{\frac{2ab}{a^2+b^2}}.
\]
\textbf{Особый случай:} когда $H=M$ (равнобедренный прямоугольный треугольник), высота и медиана совпадают и угол между ними равен $0^\circ$; формула сохраняет смысл через $\cos 0^\circ=1$.

\ex{Пример 4. Даны катеты $1$ и $2$}
При $AB=1$ и $BC=2$ имеем $AC=\sqrt5$, $BH=2/\sqrt5$, $BM=\sqrt5/2$. Следовательно,
\[
\cos\angle HBM=\frac{2/\sqrt5}{\sqrt5/2}
=\boxed{\frac45}.
\]

\section{6. Окружность на катете: полезное достроение}
Пусть $\angle BAC=90^\circ$, а окружность с диаметром $AB$ повторно пересекает гипотенузу $BC$ в точке $T$. По теореме Фалеса $\angle ATB=90^\circ$, поэтому $AT$ --- высота к гипотенузе. Если $M$ --- середина $BC$, то $AM=BM=CM=BC/2$.

\begin{center}
\begin{tikzpicture}[scale=.89]
\tkzDefPoint(0,0){C}
\tkzDefPoint(6,0){B}
\tkzDefPoint(1.5,2.598076){A}
\tkzDefPoint(1.5,0){T}
\tkzDefPoint(3,0){M}
\tkzDefMidPoint(A,B)\tkzGetPoint{O}
\tkzDrawCircle[color=accent,line width=.8pt](O,A)
\tkzDrawPolygon[line width=1pt](A,B,C)
\tkzDrawSegments[color=darkaccent,dashed](A,T A,M)
\tkzMarkRightAngle[size=.25](B,A,C)
\tkzMarkRightAngle[size=.25](A,T,B)
\tkzDrawPoints(A,B,C,T,M)
\tkzLabelPoint[above](A){$A$}
\tkzLabelPoint[below right](B){$B$}
\tkzLabelPoint[below left](C){$C$}
\tkzLabelPoint[below](T){$T$}
\tkzLabelPoint[below](M){$M$}
\end{tikzpicture}
\end{center}
\ex{Пример 5. Гипотенуза разделена в отношении $1:3$}
Если $CT:TB=1:3$, положим $CT=x$, $TB=3x$. Тогда $BC=4x$, $CM=BM=AM=2x$, а $TM=CM-CT=x$.

В прямоугольном $ATM$ катет $TM=x$ лежит напротив угла $TAM$:
\[
\sin\angle TAM=\frac{TM}{AM}=\frac{x}{2x}=\frac12.
\]
Значит, $\angle TAM=30^\circ$ и $\angle AMT=60^\circ$. Лучи $MT$ и $MB$ противоположны, поэтому $\angle AMB=120^\circ$. Равнобедренный $ABM$ даёт
\[
\angle ABC=\angle ABM=\frac{180^\circ-120^\circ}{2}
=30^\circ,
\qquad\boxed{\angle ACB=60^\circ}.
\]
\clearpage

\section{7. Что проверять при каждом решении}
\begin{enumerate}[label=\textcolor{darkaccent}{\arabic*.},leftmargin=*,itemsep=4pt]
\item \textbf{Основание для прямого угла.} Прямой угол следует из условия, теоремы Фалеса или обратного свойства медианы. На чертеже его нужно обосновать.
\item \textbf{Обозначения.} Записывайте $\angle ABC$, а при числовом значении угла указывайте знак $^\circ$ (например, $30^\circ$).
\item \textbf{Равные отрезки.} Для медианы к гипотенузе сразу отмечайте три равных отрезка от её середины к вершинам.
\item \textbf{Формульные данные.} Периметр раскрывайте как сумму сторон; площадь записывайте через удобные основание и высоту.
\item \textbf{Проверка.} Контролируйте положительность длин, теорему Пифагора, сумму углов $180^\circ$ и исходные условия.
\end{enumerate}
\textbf{Маленькая памятка.} В прямоугольном треугольнике самые полезные действия: найти гипотенузу, провести медиану к ней, провести высоту, воспользоваться равенством площадей, изучить появившиеся равнобедренные и прямоугольные треугольники.

\clearpage
\thispagestyle{fancy}
\begin{center}
{\Large\color{darkaccent}Алгоритм: поиск решения в прямоугольном треугольнике}
\end{center}
\vspace{9mm}
\tikzset{
 fstart/.style={ellipse,draw=darkaccent,thick,align=center,minimum width=42mm,minimum height=10mm,fill=white},
 faction/.style={rounded corners=2mm,draw=accent,thick,align=center,text width=72mm,inner sep=8pt,fill=white},
 fchoice/.style={diamond,aspect=2.5,draw=darkaccent,thick,align=center,text width=31mm,inner sep=2pt,fill=white},
 fside/.style={rounded corners=2mm,draw=accent,thick,align=center,text width=46mm,minimum height=19mm,inner sep=7pt,fill=white},
 farrow/.style={-{Stealth[length=3mm]},draw=darkaccent,thick}
}
\begin{center}
\begin{tikzpicture}[node distance=12mm]
\node[fstart] (s) {Начало};
\node[faction,below=of s] (angle) {Установите, где прямой угол: по условию, теореме Фалеса или обратному свойству медианы};
\node[faction,below=of angle] (mark) {Обозначьте гипотенузу $c$ и её середину $M$:\\ $AM=BM=CM=c/2$};
\node[fchoice,below=13mm of mark] (q) {Что дано в задаче?};
\node[fside,below left=17mm and 11mm of q] (len) {Стороны, периметры\\ $\Downarrow$\\ Пифагор, система уравнений};
\node[fside,below right=17mm and 11mm of q] (geo) {Высота, углы\\ $\Downarrow$\\ Две площади, синусы и косинусы};
\node[faction,below=40mm of q] (check) {Проверьте положительность длин, углы и исходные данные};
\node[fstart,below=10mm of check] (end) {Ответ};
\draw[farrow] (s)--(angle);
\draw[farrow] (angle)--(mark);
\draw[farrow] (mark)--(q);
\draw[farrow] (q.south west)--node[above left,font=\small]{длины}(len.north);
\draw[farrow] (q.south east)--node[above right,font=\small]{углы}(geo.north);
\draw[farrow] (len.south) |- (check.west);
\draw[farrow] (geo.south) |- (check.east);
\draw[farrow] (check)--(end);
\end{tikzpicture}
\end{center}
\clearpage

\section*{Тренировочный блок}
\addcontentsline{toc}{section}{Тренировочный блок}
\textbf{Путь А.} Десять заданий по возрастанию сложности. Сделайте чертёж, укажите обоснование каждого обнаруженного прямого угла и записывайте градусы с символом $^\circ$.

\textcolor{darkaccent}{\textbf{Базовый уровень}}
\begin{enumerate}[label=\textbf{\arabic*.},itemsep=9pt]
\item Гипотенуза прямоугольного треугольника равна $14$. Найдите радиус описанной окружности.
\item Медиана к гипотенузе равна $6$. Найдите гипотенузу и радиус описанной окружности.
\item В прямоугольном треугольнике $ABC$ угол $C$ прямой, $CM=m$ --- медиана к $AB$, $\angle ACM=30^\circ$. Выразите катеты $AC$ и $BC$ через $m$.
\item В прямоугольном треугольнике $ABC$ угол $B$ прямой, $AB=6$, $BC=8$. Найдите длину медианы $BM$ к гипотенузе и радиус описанной окружности.
\end{enumerate}
\textcolor{darkaccent}{\textbf{Уровень выше среднего}}
\begin{enumerate}[label=\textbf{\arabic*.},start=5,itemsep=9pt]
\item В прямоугольном $ABC$ угол $B$ прямой, $AB=5$, $BC=12$. Найдите высоту $BH$ к гипотенузе $AC$.
\item В прямоугольном $ABC$ угол $B$ прямой, $AB=3$, $BC=4$. Проведены высота $BH$ и медиана $BM$ к $AC$. Найдите $\cos\angle HBM$.
\item В треугольнике $ABC$ точка $M$ --- середина $AC$. Известно, что $AM=BM=CM=7$. Докажите, что $\angle ABC=90^\circ$, и найдите $AC$.
\item В прямоугольном треугольнике $ABC$ угол $B$ прямой, $M$ --- середина гипотенузы $AC$. Периметры треугольников $ABM$ и $BMC$ равны соответственно $18$ и $25$. Найдите все стороны $ABC$.
\end{enumerate}
\textcolor{darkaccent}{\textbf{Сложный уровень}}
\begin{enumerate}[label=\textbf{\arabic*.},start=9,itemsep=9pt]
\item В прямоугольном треугольнике $ABC$ угол $A$ прямой. Окружность с диаметром $AB$ пересекает гипотенузу $BC$ в точках $B$ и $T$. Известно, что $CT:TB=3:1$. Найдите острые углы $ABC$.
\item В прямоугольном треугольнике проведены медиана и высота из вершины прямого угла к гипотенузе. Угол между ними равен $60^\circ$. Найдите оба острых угла треугольника.
\end{enumerate}
\clearpage

\section*{Ответы}
\vspace{4mm}
\begin{center}
\renewcommand{\arraystretch}{1.65}
\begin{tabularx}{\linewidth}{@{}>{\bfseries}p{18mm}X@{}}
\toprule
№ & Ответ\\
\midrule
1 & $R=7$.\\
2 & Гипотенуза $12$, радиус $6$.\\
3 & $AC=m\sqrt{3}$, $BC=m$.\\
4 & $BM=5$, $R=5$.\\
5 & $BH=\dfrac{60}{13}$.\\
6 & $\cos\angle HBM=\dfrac{24}{25}$.\\
7 & $\angle ABC=90^\circ$, $AC=14$.\\
8 & $AB=5$, $BC=12$, $AC=13$.\\
9 & $\angle ABC=60^\circ$, $\angle ACB=30^\circ$.\\
10 & $15^\circ$ и $75^\circ$ (в любом порядке).\\
\bottomrule
\end{tabularx}
\end{center}
\vspace{12mm}
\textcolor{darkaccent}{\textbf{Итоговая памятка}}
\begin{center}
\renewcommand{\arraystretch}{1.36}
\begin{tabularx}{\linewidth}{@{}p{55mm}X@{}}
\toprule
Что замечаете & Что делаете\\
\midrule
Прямой угол и гипотенуза & Проведите медиану; отметьте равные половины гипотенузы.\\
Медиана равна полустороне & Примените обратное свойство: угол напротив этой стороны прямой.\\
Два периметра & Раскройте суммы сторон, добавьте уравнение Пифагора.\\
Высота к гипотенузе & Приравняйте площади: $ab/2=ch/2$.\\
Окружность на диаметре & Найдите угол, опирающийся на диаметр, и полезный перпендикуляр.\\
Неизвестный острый угол & Ищите равнобедренные и прямоугольные треугольники; применяйте тригонометрию.\\
\bottomrule
\end{tabularx}
\end{center}
\vfill
\begin{center}\small Лёвин Артём Александрович --- эксклюзивно для Тимофея Васильченко\end{center}
\end{document}
